\documentclass[11pt]{article}
\usepackage[margin=1.1in]{geometry}
\usepackage{amsmath,amssymb,booktabs,graphicx,hyperref}
\hypersetup{colorlinks=true,linkcolor=blue!50!black,citecolor=blue!50!black,urlcolor=blue!50!black}

\title{Measured geometry and horizon dynamics of the SU(2) gauge orbit space:\\
scale-free curvature, a localization--coherence crossover,\\
and a coherence-governed coupling law}

\author{Nitin\thanks{Independent researcher; hoo.am.aai@gmail.com. This research was conducted in close collaboration with Claude (Anthropic), which performed analysis, code, and drafting under the author's direction; per arXiv authorship policy an AI system is credited here rather than listed as an author, since authorship implies accountability only a human can hold.}}
\date{August 2026 --- preliminary; single lattice spacing per coupling, SU(2), $L\le 18$}

\begin{document}
\maketitle

\begin{abstract}
We report the first direct Monte Carlo measurements of the Riemannian geometry of the
SU(2) gauge orbit space $\mathcal{A}/\mathcal{G}$ on lattice ensembles, together with a
study of the dynamics of the approach to the Gribov horizon. Our results are fourfold.
(i)~The curvature sector is scale-free: with an empirically measured scale-setter and a
positive control, the effective mass dimensions of $K\!\cdot\!V$ and $\mathrm{Ric}_{\min}$
are bounded by $|n|\lesssim 0.02$--$0.03$ where a mass-generating quantity requires
$n=1$; structurally, a Gram--Schur inequality $\Delta \succeq M^{\dagger}L^{-1}M$ shows the
metric cannot degenerate at the horizon, and on gauge-fixed configurations the curvature is
uncorrelated with horizon proximity ($r=0.009$ while $\lambda_{\min}(M)$ varies by $167\%$).
This closes, with a mechanism, the route proposed by Singer in which orbit-space curvature
produces the mass gap. (ii)~The curvature sector's characteristic numbers are
\emph{disorder constants}: $K\!\cdot\!V$ and the trace ratio
$\mathrm{tr}\,\mathrm{Ric}/\mathrm{tr}\,T$ measured on Haar-random links reproduce the
interacting values to better than $1\%$ ($0.0794$ and $0.81576(2)$, the latter
volume-stable to $2\times10^{-5}$), with thermalized ensembles adding only per-mille
ordering corrections. (iii)~The lowest Faddeev--Popov eigenmode undergoes a
coupling-driven localization--coherence crossover ($\mathrm{IPR}\!\cdot\!V:8.7\to1.2$;
lowest-momentum weight $f_{\mathrm{IR}}:0.36\to0.95$), which persists at fixed physical
volume and is therefore not a finite-size effect; the crossover boundary is a sloped curve
in the (coupling, box-size) plane. (iv)~The per-configuration correlation between horizon
proximity and infrared gluon power --- an observable we could not locate in the literature
--- obeys a single-variable law: it is absent for $f_{\mathrm{IR}}\lesssim0.5$ and locked
at $r=-0.51..-0.53$ above $f_{\mathrm{IR}}\gtrsim0.6$, verified by a pre-registered
out-of-sample prediction in which the correlation, strong at $(\beta,L)=(2.4,10)$, was
predicted and then measured to vanish at $(2.4,16)$ ($f_{\mathrm{IR}}$ predicted $0.444$,
measured $0.443(16)$; $r$ measured $+0.02$). We also measure the dynamical approach to the
horizon ($\lambda_{\min}\sim L^{-3.08(25)}$ at $\beta=2.4$ versus the kinematic
$L^{-1.90}$) and bracket a braking size $\ell^{*}\approx 7\pm2$ in units of
$1/\sqrt{\sigma}$ at $\beta=2.2$. All claims were pre-registered before data taking; the
complete falsifier ledger, thirteen self-caught instrument bugs, and the validation suite
are part of the record. No claim herein bears on the existence of the mass gap at infinite
volume.
\end{abstract}

\section{Introduction}
The suggestion that the geometry of the gauge orbit space
$\mathcal{A}/\mathcal{G}$ might underlie the Yang--Mills mass gap goes back to
Singer~\cite{Singer78,Singer81} and the explicit curvature computations of Babelon and
Viallet~\cite{BabelonViallet}; the modern analytic form of the proposal, including the
Bakry--\'Emery substitution, appears in~\cite{Orland2018}. In parallel, the
Gribov--Zwanziger programme~\cite{Gribov,Zwanziger,VandersickelZwanziger} locates the
origin of the infrared scale at the boundary of the fundamental modular region --- the
horizon --- where the Faddeev--Popov (FP) operator loses positivity. Lattice studies of
the FP spectrum~\cite{Sternbeck05,CucchieriMendes08} and of Coulomb-gauge eigenmode
densities~\cite{GOZ} established that typical configurations approach the horizon with
growing volume.

To our knowledge, however, the \emph{geometry} side of this story has never been measured:
no curvature of $\mathcal{A}/\mathcal{G}$ has previously been evaluated on sampled
configurations, and ensemble cross-correlations between geometric, spectral, and
propagator observables --- measured jointly, configuration by configuration --- appear to
be new. This paper supplies those measurements for SU(2), and reports what they revealed.

Our methodological stance, adopted after early failures documented in the ledger, is that
every quantitative claim is stated with a falsifier fixed \emph{before} the corresponding
run, and every instrument is validated against an independent reference before use
(Sec.~\ref{sec:method}). Two headline results below (the fixed-volume crossover control and
the out-of-sample prediction of the coupling law) exist because referee-grade objections
were identified in advance and tested rather than argued.

\section{Setup and validation}\label{sec:method}
We use SU(2) Wilson lattices ($L^4$, $L=3$--$18$; $\beta=2.0$--$3.2$), a colour-class
Metropolis sampler (staples refreshed per class; detailed balance verified by
hits-independence and by agreement with an independent Creutz heat bath to $0.3\sigma$;
strong-coupling limit reproduced), and Landau gauge fixing to
$\theta<10^{-9}$. The FP operator is realized as the finite-difference Hessian of the
gauge functional (symmetric by construction; gate-checked: PSD with exactly three
colour zero modes; free-field eigenvalue exact to $10^{-10}$). Low eigenvalues are
obtained with exact deflation of the colour zero modes; we document (Appendix) a silent
failure mode of undeflated Lanczos at small $\lambda$ that biases $\lambda_{\min}$ upward
in $\sim25\%$ of solves, relevant to any similar study.

Scale setting is empirical: $d\ln a/d\beta$ is measured from Creutz ratios, with the
loop-size dependence reported explicitly ($\chi(2,2)\to\chi(4,4)$ divisors
$-0.60\to-1.9$, marching toward the asymptotic one-loop $-2.69$); dimension bounds below
use the (conservative) $\chi(3,3)$ value $-1.17$, so quoted $|n|$ are upper bounds. A
positive control (the string tension pushed through the identical pipeline) confirms
detectability of a genuine dimension.

\section{The curvature sector is scale-free, and provably blind to the horizon}
\label{sec:curv}
With the O'Neill/Babelon--Viallet form $K=3\langle\rho,\Delta^{-1}\rho\rangle$ (verified
against an independent dense construction to $9\times10^{-9}$) we measure, on corrected
ensembles with the empirical divisor,
\[
n[K\!\cdot\!V] = -0.0102(13),\qquad n[\mathrm{Ric}_{\min}] = +0.0239(70),
\]
against the $n=1$ required of a mass-generating quantity. Structurally, from the Gram
matrix of $(\partial,D)$ one obtains $\Delta \succeq M^{\dagger}L^{-1}M$: the metric
Laplacian $\Delta=D^{\dagger}D$ remains strictly positive where the FP operator $M$
crosses zero, so no quantity built from $\Delta$ and the horizontal projector alone can
detect the horizon. On $24$ Landau-fixed configurations this blindness is confirmed
directly: $\lambda_{\min}(M)$ varies by $167\%$ across the ensemble while
$\mathrm{Ric}_{\min}$ varies by $1.8\%$; $r(\lambda_{\min},\mathrm{Ric}_{\min})=+0.009$.
A configuration $12\times$ nearer the horizon than average has median curvature.

Two further facts characterize the sector. First, its numbers are \emph{disorder
constants}: on i.i.d.\ Haar links, $K\!\cdot\!V=0.0794$ and
$\mathrm{tr}\,\mathrm{Ric}/\mathrm{tr}\,T=0.81576(2)$ (volume-stable to
$2\times10^{-5}$), while thermalized ensembles differ from these by less than $1\%$
across $\beta=2.0$--$2.8$. The exact rank identity
$\overline{\mathrm{diag}\,P_H}=3/4$ makes the second constant a measure of the $G$--$H$
correlation intrinsic to the Haar measure itself; both constants are natural targets for
Weingarten-calculus evaluation. Second, $K$ is discontinuous at the trivial connection
(the zero modes of $\Delta$ at $U=1$ lift as $\epsilon^2$ nearby), so the free field is a
singular point of the geometry rather than its limit; the interacting
$K\!\cdot\!V$ is volume-flat ($0.95\%$ over $16\times$ volume) where the free-field
closed form (Appendix) rises by $7.4\%$.

Together with Sec.~\ref{sec:law} this yields the paper's negative theorem: \emph{the
curvature route to the mass gap closes because the scale is generated at the horizon
\cite{Zwanziger} and the curvature sector is structurally and measurably blind to it.}

\section{Horizon dynamics: attraction and braking}\label{sec:dyn}
With the deflated solver, the lowest FP eigenvalue falls with volume as
\[
\lambda_{\min}\sim L^{-3.08(25)}\ (\beta=2.4),\qquad
\lambda_{\min}\sim L^{-2.31(21)}\ (\beta=2.2),
\]
against the kinematic $L^{-1.90}$: at $\beta=2.4$ the ensemble is \emph{driven} to the
horizon $4.8\sigma$ faster than free momentum, with the Gribov-copy systematic bounded at
$0.16\sigma$ (best-of-five copies). Local exponents over sliding windows at $\beta=2.2$
decrease through the kinematic value near $\ell^{*}\equiv L a\sqrt{\sigma}\approx7\pm2$,
consistent with the much larger volumes of Ref.~\cite{CucchieriMendes08} where the
approach is slower than free; the existence of a turning size is thus supported at
$>3\sigma$ jointly with the literature anchor, while its dependence on $\beta$ does not
reduce to any single variable we tested and is left open. A natural constructive reading
of $\ell^*$ --- an infrared momentum cap $p^{*}=2\pi/\ell^{*}\approx0.9\sqrt{\sigma}
\approx 400\,$MeV --- is offered as an estimate only.

\section{The localization--coherence crossover}\label{sec:cross}
The lowest FP eigenmode changes character with the coupling. At $L=6$:
$\mathrm{IPR}\!\cdot\!V = 8.7(2.8)\to1.22(4)$ and lowest-momentum weight
$f_{\mathrm{IR}}=0.357(23)\to0.951(3)$ across $\beta=2.0\to3.2$: an Anderson-localized
mode reorganizes into a box-spanning, $95\%$-pure infrared wave. Because a fixed-$L$
$\beta$-scan confounds coupling with physical volume --- precisely the systematics of the
rare localized modes reported in Ref.~\cite{Sternbeck05} --- we repeated the diagnostics
at fixed physical size $\ell\approx3.1$--$3.5$: $f_{\mathrm{IR}}$ rises monotonically
$0.50(2)\to0.93(1)$, so the crossover is coupling-driven. Its boundary is a sloped curve
in $(\beta,\ell)$: near the boundary, growing the box erodes coherence
($f_{\mathrm{IR}}:0.69\to0.41$ at $\beta=2.4$ over $L=8\to18$), while deep in the
coherent phase it does not ($0.87$, flat, at $\beta=2.6$). We could not locate prior
eigenvector-shape studies of the Landau-gauge FP operator; the covariant-Laplacian
localization literature~\cite{Greensite05} reports qualitatively different behaviour
(localization persisting at weak coupling with fixed physical size).

\section{A coherence-governed coupling law}\label{sec:law}
Define, per configuration, the horizon proximity $\lambda_{\min}(M)$ and the infrared
gluon power $D(p_{\min})$ (lowest lattice momentum, averaged over orientations), and the
ensemble correlation $r$ with the plaquette partialled out. Measured across five runs and
both control directions:

\begin{center}
\begin{tabular}{lcccc}
\toprule
point & $f_{\mathrm{IR}}$ & partial $r$ & $n$ & role\\
\midrule
$\beta=2.2,\ L=10$ & $0.383$ & $+0.06$ & 48 & coupling knob (off)\\
$\beta=2.4,\ L=10$ & $0.622$ & $-0.44$ & 144 & discovery\\
$\beta=2.4,\ L=10$ (fc/bc) & $0.622$ & $-0.50/-0.55$ & 72 & replication; copy-robust\\
$\beta=2.8,\ L=10$ & $\sim0.93$ & $-0.51$ & 48 & coupling knob (on)\\
$\beta=2.6,\ L=10$ & $0.898(4)$ & $-0.534$ & 40 & \textbf{pre-registered P-ON}\\
$\beta=2.4,\ L=16$ & $0.443(16)$ & $+0.020$ & 48 & \textbf{pre-registered P-OFF}\\
\bottomrule
\end{tabular}
\end{center}

The final row is the paper's sharpest result. From the $f_{\mathrm{IR}}(\beta,L)$ grid we
predicted, in advance, that at $(\beta,L)=(2.4,16)$ the coherence would sit at $0.444$,
below the activation threshold, and that the correlation --- strong at $L=10$ at the
\emph{same coupling} --- would therefore be absent. Both predictions were confirmed
($f_{\mathrm{IR}}$ to $0.3\%$; $r=+0.02$, $0.1\sigma$ from zero). The data collapse onto
a single activation curve in $f_{\mathrm{IR}}$: off below $\approx0.5$, locked at
$r=-0.51..-0.53$ above $\approx0.6$, irrespective of whether coherence is set by the
coupling or destroyed by the volume. The sign of the correlation is anticipated at
Gribov-copy level by Ref.~\cite{SternbeckMP12}; the fixed-gauge cross-configuration
observable, and its single-variable law, appear to be new. The crossover of
Sec.~\ref{sec:cross} is the $f_{\mathrm{IR}}\approx0.5$ level set of this law.

\section{What these results do not show}
No scale is derived here. All internal relations of the theory yield dimensionless
ratios (as they must, with a single scale), and nothing in this paper bears on
$E_{\min}>0$ at infinite volume. The measurements are SU(2), at one lattice spacing per
coupling, in boxes up to $\sim L a\sqrt{\sigma}\approx10$, with no independent
reproduction yet; the coupling law's continuum fate and its SU(3) counterpart are open.

\section*{Methods, code and record availability}
All measurement scripts, GPU notebooks, the validation suite, and the complete
append-only research ledger (every claim's pre-registered falsifier, every correction
dated in place) are available at \url{https://github.com/hooamaai-create/su2-orbit-space} and permanently archived at \href{https://doi.org/10.5281/zenodo.21993867}{doi:10.5281/zenodo.21993867}; a reading guide maps the entries that
carry each of the paper's claims. We report the multiplicity honestly: of fourteen
mechanisms and patterns formulated during this work, thirteen were refuted by their own
pre-registered tests; the coherence law of Sec.~
ef{sec:law} is the survivor, and its
out-of-sample prediction was made \emph{after} that accounting.

\section*{Acknowledgements}
Analysis, code, and drafting were carried out with substantial assistance from an AI
system (Claude, Anthropic), under continuous human direction; all claims were verified
against the primary numerical record. [Funding/thanks.]

\begin{thebibliography}{99}
\bibitem{Gribov} V.~N.~Gribov, Nucl.\ Phys.\ B\,139 (1978) 1.
\bibitem{Singer78} I.~M.~Singer, Commun.\ Math.\ Phys.\ 60 (1978) 7.
\bibitem{Singer81} I.~M.~Singer, Physica Scripta 24 (1981) 817.
\bibitem{BabelonViallet} O.~Babelon, C.~M.~Viallet, Commun.\ Math.\ Phys.\ 81 (1981) 515.
\bibitem{Orland2018} P.~Orland, arXiv:1809.06318.
\bibitem{Zwanziger} D.~Zwanziger, Nucl.\ Phys.\ B\,323 (1989) 513.
\bibitem{VandersickelZwanziger} N.~Vandersickel, D.~Zwanziger, Phys.\ Rept.\ 520 (2012) 175 [arXiv:1202.1491].
\bibitem{Sternbeck05} A.~Sternbeck, E.-M.~Ilgenfritz, M.~M\"uller-Preussker, Phys.\ Rev.\ D\,73 (2006) 014502 [hep-lat/0510109].
\bibitem{CucchieriMendes08} A.~Cucchieri, T.~Mendes, Phys.\ Rev.\ D\,78 (2008) 094503 [arXiv:0804.2371].
\bibitem{GOZ} J.~Greensite, \v S.~Olejn\'ik, D.~Zwanziger, JHEP 0505 (2005) 070 [hep-lat/0407032].
\bibitem{Greensite05} J.~Greensite et al., Phys.\ Rev.\ D\,71 (2005) 114507 [hep-lat/0504008].
\bibitem{SternbeckMP12} A.~Sternbeck, M.~M\"uller-Preussker, Phys.\ Lett.\ B\,726 (2013) 396 [arXiv:1211.3057].
\bibitem{ChristLee} N.~H.~Christ, T.~D.~Lee, Phys.\ Rev.\ D\,22 (1980) 939.
\bibitem{Teper98} M.~Teper, Phys.\ Rev.\ D\,59 (1999) 014512 [hep-lat/9804008].
\end{thebibliography}

\end{document}
